\documentclass[11pt,a4paper]{article}
\usepackage[UTF8]{ctex}
\usepackage{amsmath,amssymb,amsthm,mathtools}
\usepackage{bm}
\usepackage{graphicx}
\usepackage[a4paper,left=2.5cm,right=2.5cm,top=2.6cm,bottom=2.7cm]{geometry}
\usepackage[dvipsnames]{xcolor}
\usepackage{enumitem}
\usepackage{booktabs}
\usepackage{array}
\usepackage{multirow}
\usepackage{longtable}
\usepackage{fancyhdr}
\usepackage{caption}
\usepackage{tcolorbox}
\usepackage[hidelinks,unicode]{hyperref}
\captionsetup{font=small,labelsep=quad}
\setlength{\emergencystretch}{3em}
\newenvironment{exercise}[1]{\par\medskip\noindent\textbf{习题 #1.}\itshape\ }{\par\medskip}
\newtcolorbox{keybox}[1]{colback=black!5, colframe=black!55, colbacktitle=black!75, coltitle=white,
  fonttitle=\bfseries\small, title={#1}, boxrule=0.6pt, arc=2pt,
  left=6pt, right=6pt, top=4pt, bottom=4pt, toptitle=2pt, bottomtitle=2pt, fontupper=\small}
\newcommand{\para}[1]{\par\medskip\noindent\textbf{#1}\quad\ignorespaces}
\newcommand{\figtitle}[2]{\caption*{\small\textbf{图 #1：}#2}}
\newcommand{\tabtitle}[2]{\caption*{\small\textbf{表 #1：}#2}}
\newcommand{\R}{\mathbb{R}}\newcommand{\C}{\mathbb{C}}\newcommand{\cN}{\mathcal{N}}
\newcommand{\cX}{\mathcal{X}}\newcommand{\fg}{\mathfrak{g}}\newcommand{\dd}{\mathrm{d}}
\begin{document}
\section{引言与纲要}
这是一段中文测试，包含公式 $z_i \in \C$ 与 $\cN=(2,0)$ 超共形场论。
\begin{keybox}{6d 视角下的 4d $\cN=2$ class S 理论}
\begin{itemize}[leftmargin=1.4em,itemsep=1pt,topsep=1pt]
\item 带扭曲地在穿孔黎曼面 $C$ 上约化 6d 理论 $\cX(\fg)$，得到 4d $\cN=2$ 理论。
\item Coulomb 分支由微分 $\phi_k = u_k \dd z^k$ 参数化。
\end{itemize}
\end{keybox}
\para{驯服穿刺。} 穿刺点 $z_i$ 处的一阶极点给出质量参数 $m_i$。
\begin{equation}
  \langle P_k(\Phi_z)\rangle = \frac{P_k(m_i)}{(z-z_i)^k} + \cdots, \tag{1.5}
\end{equation}
正文引用式 (1.5)，脚注测试\footnote{这是一个脚注，含 $m_i$ 与希腊字母 $\zeta(0)=-\tfrac12$。}
\begin{figure}[htbp]
\centering
\includegraphics[width=0.75\textwidth]{../figures/fig1.png}
\figtitle{1}{$\fsu(N)$ class S 理论的裤子分解示意。}
\end{figure}
\begin{exercise}{2.2}
验证沿 $x^0,x^1,\dots,x^5$ 方向放置的平直 M5 膜保持超对称性。
\end{exercise}
\begin{table}[htbp]
\centering
\begin{tabular}{llll}
\toprule
超代数 & 共形代数 & R 对称性 & $Q$ 与 $S$ \\
\midrule
3d $\cN\le 8$ & $\fsp(4,\R)$ & $\fso(\cN)$ & $(4,\cN)$ \\
\bottomrule
\end{tabular}
\tabtitle{2}{Nahm 分类（测试用）。}
\end{table}
\end{document}
